In this work, we investigated two directions for improving transformer-based cell tracking: efficient sparse attention and self-supervised pretraining.

\subsection{Efficient Sparse Attention}

Our comprehensive benchmarking reveals a clear picture. At the cell counts spanned by the datasets analysed in this work, vanvliet~\cite{vanvliet2018} and DeepCell~\cite{schwartz2023}, whose per-sample $N$ distributions we characterize statistically in Section~\ref{sec:dataset_stats}, the attention mechanism is a negligible fraction of total training time ($1.62$--$1.64$\,min/epoch, identical across all $K$). The corrected benchmark (using CachedDistAttention and including the $\mathcal{O}(N^2)$ KNN index cost in timing) shows that dense masked attention is $1.7\times$ faster than mask-KNN sparse attention at $N \approx 166$, so sparse attention is strictly slower at the cell-tracking scale. We address the per-layer pairwise distance computation (which accounted for $82\%$ of attention time in the old per-layer baseline) with CachedDistAttention, which computes pairwise distances once and shares them across layers with identical semantics and no approximation. On the RTX A500 (fp16, $d_{\text{model}}=320$, $n_{\text{head}}=8$, batch $2$, $L=12$), its measured end-to-end speedup over the per-layer baseline is regime-dependent: $1.61\times$ at $N=128$ and $1.82\times$ at $N=256$---the small-$N$ regime that dominates the vanvliet dataset (Section~\ref{sec:dataset_stats})---but only $1.02$--$1.06\times$ for $N=512$--$2048$ and $1.20\times$ at $N=4096$, before running out of memory at $N=8192$ alongside the per-layer baseline. A blanket $\sim 2\times$ figure therefore holds only for $N \leq 256$.

For the cell-tracking scale ($N \leq 512$), we recommend dense masked attention with CachedDistAttention as the default architecture. The corrected benchmark shows that at this scale, dense masked attention ($0.199$--$0.344$\,ms at $N = 128$--$2048$) is faster than both mask-KNN ($0.335$--$0.396$\,ms) and gather-KNN ($0.400$--$0.480$\,ms), and inherits the validated tracking accuracy of the Trackastra baseline (TRA $0.9963$, AOGM $51.0$; single seed). For $N \geq 4096$, mask-KNN with $K=4$--$16$ becomes $1.5$--$1.75\times$ faster than dense masked, and gather-KNN with $K=4$--$16$ becomes $2.0$--$2.25\times$ faster---but this regime is not reached by the analysed datasets. Token reordering provides negligible benefit ($0.97$--$1.03\times$) on random coordinate data, and its value on the cell centroids of the analysed datasets remains unverified. NSA is architecturally mismatched to this domain ($29\times$ slower at $N=512$, out-of-memory at $N=8192$ for $K \geq 4$) and should not be used for bidirectional cell tracking transformers.

\subsection{Spatial Cutoff Ablation: The Cutoff Is Purely Computational}

The Trackastra paper~\cite{gallusser2024trackastra} fixes the spatial cutoff at $d_{\max} = 256$\,px but never ablates it (Section~\ref{sec:eval_cv}). Our 6-fold leave-one-condition-out cross-validation is the first proper test of whether the spatial cutoff matters for tracking accuracy. The results are unambiguous: removing the spatial cutoff (DenseFlashAttention, FlashAttention-2 backend) produces identical tracking accuracy to the masked baseline ($\Delta$TRA $= -0.0003$, $\Delta$AOGM $= +4.0$ errors out of $\sim 225$). The spatial cutoff is therefore \emph{purely computational}---it exists to reduce the effective attention density at high $N$, not to regularize the model or improve tracking accuracy.

This finding has a direct architectural implication: at vanvliet scale (median $N = 62$, mean $N = 165.7$), the optimal architecture is \textbf{DenseFlashAttention} (no spatial cutoff mask, FlashAttention-2 backend), which yields a $1.78\times$ attention kernel speedup with zero accuracy cost.

\subsection{Amdahl's Law: Attention Is Not the Training Bottleneck}

Despite the $1.78\times$ attention kernel speedup, the measured training speedup is only $\sim 3\%$ (Section~\ref{sec:eval_amdahl}). This is because attention occupies only $\sim 5\%$ of total step time at $N \approx 166$, where data loading from Lustre ($\sim 40$--$60$\,ms), data augmentation ($\sim 15$\,ms), feed-forward layers ($\sim 20$\,ms), and the backward pass ($\sim 30$\,ms) dominate. Amdahl's Law predicts a $2.2\%$ training speedup from a $1.78\times$ attention speedup at $5\%$ attention share, which matches the measured $\sim 3\%$ from the 6-fold CV.

The implication is that \emph{attention is not the training bottleneck} at vanvliet scale. The bottleneck is data loading and feature extraction, which run on CPU and are largely invisible to GPU-side attention optimizations. This is consistent with the profiler analysis (Section~\ref{sec:app_profiler}), which shows that attention occupies only a small fraction of total step time at $N \approx 166$.

\subsection{Other Optimizations Provide Negligible Training Speedup}

The optimized variant (cached-dist-attn branch, with CachedDistAttention, fast-regionprops, and vectorized blockwise\_norm) trains at $3.7$\,min/epoch, which is \emph{not faster} than the vanilla baseline ($3.6$\,min/epoch; Table~\ref{tab:app_per_epoch}). The individual optimizations have the following isolated speedups:

\begin{itemize}[noitemsep]
    \item \textbf{FlashAttention:} $1.78\times$ on attention kernel, but attention is only $\sim 5\%$ of step time $\Rightarrow$ $2.2\%$ training speedup.
    \item \textbf{Vectorized blockwise\_norm:} $1.96\times$ on the normalization component, but normalization is only $\sim 15\%$ of step time $\Rightarrow$ $\sim 6\%$ training speedup (theoretical).
    \item \textbf{CachedDistAttention:} $1.5\times$ on cdist computation, but cdist is only $\sim 3\%$ of step time $\Rightarrow$ $\sim 1\%$ training speedup.
    \item \textbf{fast-regionprops:} $10$--$80\times$ on CPU feature extraction, but this runs in the data loader parallel with GPU training $\Rightarrow$ $0\%$ training speedup.
\end{itemize}

The combined theoretical training speedup is $\sim 9.5\%$, but the measured speedup is $\sim 3\%$. The gap is because the optimizations target overlapping fractions of step time, and because fast-regionprops runs in the data loader (CPU, parallel with GPU) and therefore provides zero training speedup.

\subsection{CNN Feature Injection}

We evaluated eight CNN feature injection variants on the H100. All CNN-augmented models were outperformed by the $7$D regionprops-only baseline, with validation loss $3.8$--$25.7\times$ worse. CONCAT injection (features blended before LayerNorm) achieved the best result among CNN variants (validation loss $0.011$, $3.8\times$ worse than baseline), confirming that post-LayerNorm additive injection creates an irreversible local minimum trap that even scheduled $\lambda(t)$ decay cannot escape. Cross-dataset evaluation on DeepCell confirmed that the baseline generalizes best (TRA $0.990$, cHOTA $0.956$, AOGM $592$), with CNN variants showing small but consistent degradation. Edge probing revealed diminishing but nonzero returns from feature engineering: DINOv2 (balanced accuracy $0.896$) improves only $+0.036$ over $7$D regionprops ($0.860$), and HOCT-style $19$D features close approximately one-third of this gap with zero overhead. The key finding is that $7$D regionprops are surprisingly sufficient for tracking. Consistent with the post-LayerNorm injection trap identified in the CNN experiments, we attribute the failure of feature injection to an architectural cause rather than to feature quality; the diminishing-but-nonzero returns in edge probing suggest---though they do not prove---that architectural changes, rather than feature engineering, are the more promising path to improvement.

\subsection{SSL Pretraining}

Both identity BCE and NT-Xent contrastive pretraining on $7$-dimensional regionprops features failed to improve over a from-scratch baseline: identity BCE reached a validation loss of $0.404$ versus $0.416$ for the from-scratch baseline after fine-tuning, and NT-Xent collapsed to an inter-cell cosine similarity of $0.89$ with a downstream TRA of $65.1\%$ versus $64.8\%$ for the random baseline. The HOCT paper~\cite{bragantini2026hoct} reports findings consistent with an architectural explanation: their $19$D handcrafted features achieve state-of-the-art tracking \emph{without any visual encoder}, while their edge-centric architecture enables a $59\%$ AOGM reduction via active learning probes, which the HOCT authors report as roughly a $30\times$ larger reduction than a comparable node-centric linear probe. This is consistent with---and we attribute the dominant cause of the failure to---the architectural paradigm (node-centric vs.\ edge-centric) rather than to feature quality; however, the decisive in-architecture test (a coordinate-free variant operating on $7$D regionprops) was never run, so the attribution remains hedged. SSL for cell tracking may still benefit from richer representations, but even with DINOv2-level features the node-centric architecture saturates near $0.90$ balanced accuracy on edge prediction, suggesting a ceiling that is architectural rather than representational.

\subsection{Limitations}

We acknowledge several limitations of this work. Cross-dataset evaluation is restricted to DeepCell (12 sequences, fluorescence microscopy)~\cite{schwartz2023}; the dataset resides on the HPC cluster, so its per-sample cell-count statistics could not be verified from the local workspace (Section~\ref{sec:dataset_stats}), and additional benchmarks on other datasets would strengthen generalizability claims. The CNN feature-injection experiments, while covering eight architectural variants (Table~\ref{tab:app_cnn}), rely on a single CNN backbone (a lightweight 128-D ConvNet encoder pretrained with NT-Xent); alternative visual encoders were not tested in the injection setting. Our experiments ran on two GPU architectures: an NVIDIA RTX A500 Laptop GPU (4 GB VRAM, Ampere) for the attention benchmarks and the K-sweep, and an NVIDIA H100 (80 GB VRAM, Hopper) on the TU Dresden Capella cluster, on which the CNN feature-injection experiments and the full-scale training runs were conducted. No substantial hyperparameter optimization was performed---the hyperparameters of the Trackastra baseline were adopted and kept fixed, so optimal configurations for the sparse variants were not searched. All runs used a single random seed (seed 42), leaving inter-run variance unquantified; the consistent $\pm 0.001$ TRA spread across configurations (Section~\ref{sec:eval_cv}) suggests this effect is minor, but it is not directly measured.

Several benchmark-level caveats affect the interpretation of our throughput results. The corrected H100 benchmark (job 3920627, \texttt{--with-knn} flag) includes the KNN index cost in timing, matching training behavior; the legacy A500 benchmark excludes it. The K-sweep training runs (Section~\ref{sec:eval_cv}) exclude KNN precomputation from the per-epoch wall-clock measurement, and mask-KNN's $\mathcal{O}(N^2)$ Boolean Mask has been characterized at $N \leq 512$---its scalability at $N \geq 2048$ was not measured during the K-sweep training runs but was benchmarked separately on the H100 (job 3920627). Our maximum measured sequence length is $N = 8192$; performance at $N > 64{,}000$ (the regime where NSA~\cite{yuan2025native} claims its advantage) is extrapolated from published results rather than direct measurement. While fast-regionprops~\cite{weigert2025fastregionprops} offers a $10$--$80\times$ speedup over the standard scikit-image regionprops implementation and is integrated into the training pipeline, it was not incorporated into the attention benchmark harness. Additionally, although cHOTA was computed alongside TRA and AOGM for completeness on the DeepCell evaluation (Table~\ref{tab:app_deepcell}), our primary accuracy analysis throughout this work relies on TRA and AOGM as the measures of record.

The most fundamental limitation is architectural rather than parametric, although this conclusion is itself hedged. Our node-centric attention mechanism collapses pairwise edge representations to a single scalar, whereas HOCT's edge-centric architecture~\cite{bragantini2026hoct} preserves high-dimensional edge tokens amenable to discriminative probing: a logistic regression head on $288$D edge features achieves a $59\%$ AOGM reduction with only $400$ annotations, which the HOCT authors report as roughly a $30\times$ larger reduction than a comparable node-centric linear probe ($-2.2\%$ AOGM). We attribute this ceiling to the node-centric architecture rather than to feature quality, consistent with the post-LayerNorm injection trap and with the diminishing-but-nonzero feature returns observed in edge probing, but the decisive in-architecture test (a coordinate-free variant operating on $7$D regionprops) was never run, so a feature-only explanation cannot be fully excluded. SSL evaluation is therefore not a definitive judgment on SSL for cell tracking; the observed failure is consistent with an architectural bottleneck, but the attribution remains tentative.

\subsection{Future Work}

Our results point toward an architectural next step rather than a feature-engineering one. The post-LayerNorm injection trap, the saturation of node-centric edge probing near $0.90$ balanced accuracy, and the HOCT authors' roughly $30\times$ larger AOGM reduction from edge-centric probing all suggest that the node-centric paradigm itself is the bottleneck. The most promising direction is therefore to integrate HOCT's edge-centric paradigm into our sparse-attention framework: construct explicit high-dimensional edge tokens and attend over them as HOCT does, while retaining our mask-KNN sparse attention and CachedDistAttention, whose measured benefit is confined to the small-$N$ regime ($N \leq 256$) that dominates the datasets analysed here and does not extend to very large $N$. An edge-centric variant of our tracker would let us test directly whether the SSL and feature-injection failures persist once edge representations are preserved rather than collapsed to a single scalar, and whether active-learning probes on edge tokens achieve AOGM reductions comparable to those reported for HOCT.

Secondary directions follow from the hedged conclusions of this work. Alternative self-supervised objectives, such as BYOL, SimSiam, or DINO-style methods, may avoid the collapse observed with NT-Xent (inter-cell cosine similarity $0.89$, no downstream gain) and remain worth investigating once an edge-centric architecture is in place. Improved feature engineering is likewise not ruled out by our results: edge probing showed diminishing but nonzero returns ($+0.036$ balanced accuracy for DINOv2 over $7$D regionprops), and the decisive in-architecture test of whether richer features help within an edge-centric model was never run. A linker ablation (greedy vs.\ LAP vs.\ ILP), multi-seed training, and a proper hyperparameter study would also strengthen the statistical claims made here, which currently rest on single-seed runs.

\clearpage
